\documentclass[11pt]{article}

\usepackage[margin=1in]{geometry}
\usepackage{amsmath,amssymb,amsthm,mathtools}
\usepackage{bm}
\usepackage{enumitem}
\usepackage{booktabs}
\usepackage{hyperref}

\setlist{nosep,leftmargin=*}
\setlength{\parskip}{0.5em}
\setlength{\parindent}{0pt}

\title{The Lie Derivative in Atmospheric Dynamics\\
\large Notes, Worked Examples, and Homework}
\author{UAH Course: Geometric and Topological Methods in Atmospheric Dynamics}
\date{}

\newcommand{\Lie}{\mathcal{L}}
\newcommand{\DD}{\mathrm{D}}
\newcommand{\R}{\mathbb{R}}
\newcommand{\ii}{\hat{\mathbf{i}}}
\newcommand{\jj}{\hat{\mathbf{j}}}
\newcommand{\kk}{\hat{\mathbf{k}}}

\begin{document}
\maketitle

\section*{1. Why the Lie Derivative?}

In atmospheric dynamics, the material derivative tracks how a quantity changes along a parcel path. For geometric objects (vectors, flux forms, vortex tubes, fronts), this is incomplete by itself because the flow deforms the ambient space by stretching, shearing, and rotating.

The Lie derivative $\Lie_{\mathbf{u}}$ measures transport \emph{relative to the geometry of the flow map}. In short:
\begin{itemize}
  \item Material derivative: physical rate of change along trajectories.
  \item Lie derivative: geometric rate of change under pushforward by the flow.
\end{itemize}

\section*{2. Core Definitions and Identities}

\subsection*{2.1 Vector-field definition}

For vector fields $\mathbf{u}$ (velocity) and $\mathbf{A}$,
\[
\Lie_{\mathbf{u}}\mathbf{A} \equiv [\mathbf{u},\mathbf{A}] = (\mathbf{u}\cdot\nabla)\mathbf{A} - (\mathbf{A}\cdot\nabla)\mathbf{u}.
\]

Using
\[
\frac{\DD \mathbf{A}}{\DD t} = \partial_t \mathbf{A} + (\mathbf{u}\cdot\nabla)\mathbf{A},
\]
we get
\[
\Lie_{\mathbf{u}}\mathbf{A}
= \frac{\DD \mathbf{A}}{\DD t} - \partial_t\mathbf{A} - (\mathbf{A}\cdot\nabla)\mathbf{u},
\]
or equivalently
\[
\frac{\DD \mathbf{A}}{\DD t}
= \partial_t\mathbf{A} + \Lie_{\mathbf{u}}\mathbf{A} + (\mathbf{A}\cdot\nabla)\mathbf{u}.
\]

The term $(\mathbf{A}\cdot\nabla)\mathbf{u}$ is the stretching/tilting contribution.

\subsection*{2.2 Frozen-in criterion (important correction)}

A common misconception is ``frozen-in'' means $\DD\mathbf{A}/\DD t=0$. That is generally false in deforming flows.

\textbf{Correct criterion:}
\[
\mathbf{A}\text{ is Lie-dragged by }\mathbf{u}
\iff \Lie_{\mathbf{u}}\mathbf{A}=0.
\]

If $\Lie_{\mathbf{u}}\mathbf{A}=0$, then $\mathbf{A}$ may still change in magnitude/direction due to deformation via $(\mathbf{A}\cdot\nabla)\mathbf{u}$.

\subsection*{2.3 Differential-form form (Cartan identity)}

For any differential form $\alpha$,
\[
\Lie_{\mathbf{u}}\alpha = d(\iota_{\mathbf{u}}\alpha) + \iota_{\mathbf{u}}(d\alpha).
\]

For transport equations, the geometric advection law is
\[
\partial_t \alpha + \Lie_{\mathbf{u}}\alpha = 0.
\]

This is the coordinate-free template behind many conservation laws.

\section*{3. Atmospheric Geometric Transport}

\subsection*{3.1 Mass as a 3-form}

Define mass form
\[
\mu = \rho\,dx\wedge dy\wedge dz.
\]
Continuity is simply
\[
\partial_t \mu + \Lie_{\mathbf{u}}\mu = 0.
\]

\subsection*{3.2 Vorticity as a 2-form}

Let velocity 1-form be $v=\mathbf{u}^{\flat}$ and vorticity 2-form
\[
\Omega = dv.
\]
In inviscid barotropic flow (with conservative body force),
\[
\partial_t\Omega + \Lie_{\mathbf{u}}\Omega = 0.
\]

\subsection*{3.3 Ertel PV as a ratio of transported forms}

Let $\theta$ be potential temperature and define the 3-form identity
\[
Q\,\mu = d\theta \wedge \Omega.
\]
Under adiabatic, inviscid assumptions, both sides are Lie-transported, implying
\[
\frac{\DD Q}{\DD t}=0.
\]
This avoids long vector-identity derivations.

\section*{4. Worked Example}

\subsection*{Example: Sheared jet}

Take
\[
\mathbf{u}(x,y,z)=(U_0+Sy)\,\ii,
\qquad
\mathbf{A}=A_0\,\jj,
\]
with $U_0,S,A_0$ constants, and assume locally steady $\partial_t\mathbf{A}=0$ and $\DD\mathbf{A}/\DD t=0$.

From
\[
\Lie_{\mathbf{u}}\mathbf{A} = \frac{\DD\mathbf{A}}{\DD t}-\partial_t\mathbf{A}-(\mathbf{A}\cdot\nabla)\mathbf{u},
\]
we get
\[
\Lie_{\mathbf{u}}\mathbf{A}=-(\mathbf{A}\cdot\nabla)\mathbf{u}.
\]
Now
\[
(\mathbf{A}\cdot\nabla)=A_0\partial_y,
\quad
A_0\partial_y\big((U_0+Sy)\ii\big)=A_0S\,\ii.
\]
Hence
\[
\boxed{\Lie_{\mathbf{u}}\mathbf{A}=-A_0S\,\ii.}
\]
So the vector is not geometrically invariant; shear rotates/tilts the transported direction.

\section*{5. Homework}

\subsection*{Problem 1: Convective stretching}

Consider an idealized updraft
\[
\mathbf{u}(x,y,z)=\sigma z\,\kk,
\]
with constant $\sigma>0$, and a vertical vortex filament
\[
\mathbf{B}=B\,\kk.
\]
Assume local steadiness $\partial_t\mathbf{B}=0$ and frozen-in condition $\Lie_{\mathbf{u}}\mathbf{B}=0$.

\begin{enumerate}[label=\arabic*.]
  \item Compute $\DD\mathbf{B}/\DD t$.
  \item Solve for $B(t)$ along a parcel trajectory.
  \item Interpret physically in terms of vortex amplification.
\end{enumerate}

\subsection*{Problem 2: Uniform zonal advection of a wave-like vector field}

Let
\[
\mathbf{u}=U_0\,\ii,
\qquad
\mathbf{C}(x)=\cos(kx)\,\jj,
\]
with constants $U_0,k$.

\begin{enumerate}[label=\arabic*.]
  \item Compute $\Lie_{\mathbf{u}}\mathbf{C}=[\mathbf{u},\mathbf{C}]$ directly.
  \item Determine whether $\mathbf{C}$ is Lie-dragged by the flow.
\end{enumerate}

\section*{6. Solutions (Instructor Key)}

\subsection*{Solution 1}

From
\[
\frac{\DD\mathbf{B}}{\DD t} = \partial_t\mathbf{B} + \Lie_{\mathbf{u}}\mathbf{B} + (\mathbf{B}\cdot\nabla)\mathbf{u},
\]
with first two terms zero,
\[
\frac{\DD\mathbf{B}}{\DD t}=(\mathbf{B}\cdot\nabla)\mathbf{u}.
\]
Since $(\mathbf{B}\cdot\nabla)=B\partial_z$ and $\mathbf{u}=\sigma z\,\kk$,
\[
\frac{\DD\mathbf{B}}{\DD t}=B\sigma\,\kk.
\]
Thus scalar amplitude satisfies
\[
\frac{\DD B}{\DD t}=\sigma B
\quad\Rightarrow\quad
B(t)=B(0)e^{\sigma t}.
\]
Interpretation: vertical stretching amplifies vertical vorticity/tube intensity.

\subsection*{Solution 2}

Compute bracket terms:
\[
(\mathbf{u}\cdot\nabla)\mathbf{C}
=U_0\partial_x\big(\cos(kx)\,\jj\big)
=-kU_0\sin(kx)\,\jj,
\]
\[
(\mathbf{C}\cdot\nabla)\mathbf{u}
=\cos(kx)\partial_y(U_0\,\ii)=0.
\]
Hence
\[
\boxed{\Lie_{\mathbf{u}}\mathbf{C}=-kU_0\sin(kx)\,\jj\neq 0.}
\]
So $\mathbf{C}$ is not Lie-invariant under this flow.

\section*{7. Quick Reference Box}

\begin{tabular}{@{}ll@{}}
\toprule
Concept & Formula \\
\midrule
Vector Lie derivative & $\Lie_{\mathbf{u}}\mathbf{A}=(\mathbf{u}\cdot\nabla)\mathbf{A}-(\mathbf{A}\cdot\nabla)\mathbf{u}$ \\
Material/Lie relation & $\DD\mathbf{A}/\DD t=\partial_t\mathbf{A}+\Lie_{\mathbf{u}}\mathbf{A}+(\mathbf{A}\cdot\nabla)\mathbf{u}$ \\
Frozen-in criterion & $\Lie_{\mathbf{u}}\mathbf{A}=0$ \\
Cartan identity & $\Lie_{\mathbf{u}}\alpha=d(\iota_{\mathbf{u}}\alpha)+\iota_{\mathbf{u}}d\alpha$ \\
Advection of form & $\partial_t\alpha+\Lie_{\mathbf{u}}\alpha=0$ \\
\bottomrule
\end{tabular}

\end{document}
